\begin{claim}{E2.1}{Known}{The structure sheaves of the comb members}
Let $X_m$ be any sufficiently small smooth period section obtained by
the $m$-tooth construction of
\citet[Theorem 1.1]{ViaclovskyCombs2026}, with
$m\ge12$, fixed limiting data at the three distinguished values, and
the standard logarithmic-transformation data. Write
$f_m:X_m\to B$. Then
\[
 R^1f_{m*}\OO_{X_m}=\OO(a_m)\oplus\OO(b_m),\quad
 a_m,b_m\ge0,\quad a_m+b_m=2m-1,
\]
\[
 R^2f_{m*}\OO_{X_m}=\OO(2m-1),\qquad
 (h^{0,0},h^{0,1},h^{0,2},h^{0,3})=(1,2m+1,2m,0).      \tag{E.1}
\]
The individual summands are not specified by $m$ in the construction.
\end{claim}
\begin{proof}
The paper's equation (7.5) proves that attaching and smoothing one
twistor conic increases the determinant degree of its quotient bundle
by two. Its Proposition 9.1 identifies the initial degree as $-1$,
so the smoothed quotient has degree $2m-1$. The geometric gluing
comparison in its Section 6.3 identifies that quotient with
$R^1f_{m*}\OO_{X_m}$; the fixed limiting periods preserve the same
cyclic representations at the finite points and the same
normalization complex at the cusp. Thus its Theorem 2.3 applies to
these actual full scheme-theoretic fibres and gives the cup-product
isomorphism $\det R^1f_{m*}\OO=R^2f_{m*}\OO$.
Its Proposition 3.2, whose hypotheses are connected fibres, torus
general fibre and the topology of $S^6$, gives the two possible
splitting alternatives: both degrees nonnegative, or $(-1,0)$.
The positive determinant degree excludes the latter. We have now
derived the stated sheaves from the source comparisons in precisely
the form needed here. Leray on $B$ gives $h^{0,1}=a_m+b_m+2$,
$h^{0,2}=a_m+b_m+1$ and $h^{0,3}=0$, proving (E.1).
The theorem is used for the smoothed section, not for the singular
comb, which is not a small deformation of $X_0$.
\end{proof}

